\chapter{謝辞}
\label{chp:thanks}
某大学の優秀な女学生が卒業生代表として謝辞の冒頭でこう話した．「皆が皆，同じ経験をして，同じように感じるならば，わざわざ言葉で表現する必要はない．
見事な定型文と美辞麗句の裏側にあるのは完全な思考停止だ．」と．
その通りである．私が頑張らなければここまで来れかったはずで，私がここまでこれたのは全て私が行動したからなのだ．
私こそ，大学生活において最も感謝すべき人なのである．
支えてくれた人もいるが，残念ながら私のことを大学に対して批判的な態度であると揶揄する人もいた．
しかし，私は素晴らしい学績を納めたので「おかしい」ことを口にする権利がある．
大した仕事もせずに，自分の権利ばかり主張する人間とは違うのである．
結局，全ては私自身のおかげに他ならないのである．

と，私も堂々と言えたならばどれほどよかっただろうか．
皆が知っての通り，私はそんなに優秀な学生ではない．彼女のようにすべての年度での成績優秀者ではない，
学生の提言の優秀賞、卒業論文の最優秀賞などの素晴らしい学績を獲得したわけでもない，
ただのしがない学生に過ぎないのである．そんな私にそんな開き直りをする権利がどこにあるだろうか，いやない．
あれほど大層なことを言えるのは，それを言えるだけの実力と自信があるからである．
私は根拠のない自信を持て余しては，そこら中に撒き散らしているだけの大うつけ者なのである．
そんな大うつけ者が，こうして研究を終えることができたのも，
結局，全ては皆さまのおかげに他ならないのである．

それでは，改めて私から以下の方に向けて，懇切丁寧に謝辞の言葉を捧げる．

本研究の一番の貢献者といっても過言ではないお方，本学科教授今堀慎治先生．
タイトル決めで悩んでいたとき，素晴らしいアドバイスを下さったお方，同じく本学科助教李東珍先生．
本研究について興味を示してくださり，今後の院生としての在り方を説いてくださったお方，
インターン先の先輩のヒジカタ氏．
本研究にて，心の支えとなってくれた安西いおり先輩や伊藤じゅん先輩．
実験用の画像をこれでもかと言うほどくれた私の友人の方々
（塚本くん，神山くん，泉くん，林くん，宮尾くん，ビンセントくん，あやかさん，陳くん，高橋くん，坂本くん）．
ゼミでの進捗発表では的確な質問やアドバイスを，添削においては相談にのって下さった丹治春人くん．
Processingの表現できることやその面白さ，魅力を最大限に伝えてくれた山本晶偉くん．
本論文作成に至り，画像提供だけでなく掲載許可まで出していただき，私に大いに貢献してくれた猪口理沙さん．
後輩でありながら，その多くの知識と技術で私を献身的に支えてくれた今枝俊輔くん．
学部卒業まで最後まで支えてくれた両親．関係した全ての方に厚く御礼申し上げる．